% basic nastaveni
\documentclass[12pt]{article}
\setlength{\parindent}{0pt}
\setlength{\parskip}{1em}
\usepackage[utf8]{inputenc}
\usepackage[czech]{babel}
\usepackage[T1]{fontenc}
\usepackage{graphics}
\usepackage{caption}
\usepackage{hyperref}
\usepackage{xcolor}
\usepackage{float}

% matematicke package
\usepackage{amsmath}
\usepackage{amsfonts}
\usepackage{amssymb}
\usepackage{geometry}
\geometry{margin=2.5cm}
\usepackage{pgfplots}
\pgfplotsset{compat=1.18}
\usepackage{mathrsfs}

\title{Maturitní otázka číslo 15: Tělesa}
\author{}
\date{}

\begin{document}
\maketitle
\subsection*{Co to jsou tělesa}
Prostorově omezené souvislé útvary se nazývají tělesa. Jejich hranicí jsou uzavřené plochy. Vnitřní
body jsou body, které se nacházejí uvnitř prostoru. Hraniční body se nacházejí na plochách daného
tělesa. Každé těleso musí mít nenulový objem. Pokud všechny úsečky tvořené všemi dvojicemi bodů
tělesa leží uvnitř tělesa, jde o konvexní těleso, jinak jde o nekonvexní.
\subsection*{Mnohostěny}
Mnohostěny se skládají z vrcholů, hran a stěn. Vrcholy jsou body, kde se stýkají hrany, hrany jsou
úsečky, kde se stýkají dvě stěny a stěny jsou mnohoúhelníky tvořící povrch tělesa. Dále v nich
rozlišujeme stěnové a tělesové úhlopříčky. Stěnové úhlopříčky spojují vrcholy vrámci jedné stěny,
zatímco tělesové spojují vrcholy, které neleží v téže stěně.
\subsubsection*{Hranoly}
Hranoly jsou tělesa, které mají dvě shodné podstavy v rovnoběžných rovinách a boční stěny, které
jsou tvořeny rovnoběžníky. Hranoly rozlišujeme na n-boké podle toho, kolik hran má jejích podstava.
Rozlišujeme dva specifické hranoly a to kolmý hranol, jehož boční stěny jsou kolmé k podstavám, takže
všechny boční stěny jsou obdelníky nebo čtverce. Dalším z nich je pravidelný hranol, což je typ
kolmého hranolu, jehož podstavou je pravidelný n-úhelník. U pravidelného hranolu platí, že délka
všech tělesových úhlopříček je stejná.

$V = S_p \cdot v$, $S_p$ je obsah podstavy a $v$ je výška hranolu

$S = 2 \cdot S_p + S_{pl}$, $S_{pl}$ je obsah pláště.

Kvádr je speciální typem hranolu. Pokud je jeho podstavou obdélník, jde o kolmý čtýřboký hranol, který
má délku podstav $a, b$ a výšku $c$. Speciálním typem kvádru je čtvercový kvádr, který má jako 
podstavu čtverec je to pravidelný čtyřboký hranol. Kvádr má čtyři tělesové úhlopříčky, které jsou
stejně dlouhé a jejich délka lze vypočítat jako $u_t = \sqrt{a^2 + b^2 + c^2}$.

$V = a \cdot b \cdot c$

$S = 2(a \cdot b + b \cdot c + a \cdot c)$

Krychle je pravidelný čtyřboký hranol, který má jako podstavu čtverec délky $a$ a výšku taky $a$.
Všechny tělesové úhlopříčky mají stejnou délku $u_t = a \cdot \sqrt{3}$ a všechny stěnové $u_s =
a \cdot \sqrt{2}$.

$V = a^3$

$S = 6 \cdot a^2$

\subsubsection*{Jehlan, komolý jehlan}
Jehlan je mnohosťen, který má podstavu jako libovolný n-úhelník a plášť tvořený trojúhelníky, které
se sbíhají ve vrcholu značeném $V$. Speciálním typem jehlanu je pravidelný jehlan, který má jako
podstavu pravidelný n-úhelník.

$V = \frac{1}{3} \cdot S_p \cdot v$

$S = S_p + S_{pl}$

Komolý jehlan vznikne tím, že uřízneme vršek jehlanu rovinou, která je rovnoběžná s rovinou podstavy.
Komolý jehlan má tedy dvě podstavy $S_1, S_2$, které jsou si podobné.

$V = \frac{1}{3} \cdot v \cdot (S_1 + \sqrt{S_1 \cdot S_2} + S_2)$

$S = S_1 + S_2 + S_{pl}$
\subsubsection*{Pravidelné mnohostěny}
Pravidelné mnohostěny, někdy také nazývané Platónská tělesa je specialní skupina konvexních 
mnohostěnů. Tyto mnohostěny jsou tvořeny stejnými n-úhelníky a v každém vrcholu se stýká stejný
počet hran. Patřá mezi ně pravidelný čtyřstěn, šestistěn(krychle), osmistěn, dvanáctistěn a
dvacetistěn. Pro Platónská tělesa platí tyto dva principy. Eulerova věta o mnohostěnech říká,
že pro každý konvexní mnohostěn, tedy i platónská těles, platí vztah $V - H + S = 2$. $V$ značí
počet vrcholů, $H$ počet hran a $S$ počet stěn. Dále u Platónských těles platí princip duality.
Mějme nějaké platósnké těleso. Pokud spojíme středy stěn tohoto tělesa, získáme jiné platónské
těleso. Krychle a osmistěn jsou mezi sebou duální a také dvanáctistěn a dvacetistěn jsou duální
dvojicí. Čtyřstěn je samoduální, tedy po spojení středů stran vznikne další čtyřstěn.

\subsection*{Rotační tělesa}
Rotační těleso je těleso, které vznikne rotací rovinného útvaru kolem přímky, která se nazývá osa 
rotace.
\subsubsection*{Válec}
Rotační válec vznikne rotací obdélníků kolem jedné z jeho stran. Strana kolem které se rotuje udává
výšku $v$ a druhá tvoří poloměr $r$ podstavy. Podstavy válce jsou dva shodné kruhy, které leží v
rovnoběžných rovinách. Speciálním válcem je rovnostranný válec, který má výšku rovnu průměru válce a
má osový řez, kterým je čtverec.

$V = S_p \cdot v = \pi r^2 v$

$S = 2 \cdot S_p + S_{pl} = 2 \pi r^2 + 2 \pi r v$

\subsubsection*{Kužel, komolý kužel}
Kužel vzniká rotací pravoúhlého trojúhelníku kolem jedné z jeho odvěsen. Odvěsna ležicí na ose
rotace udává výšku $v$, druhá odvěsna udává poloměr $r$ a přepona stranu $s$, přičemž pořád platí
$s^2 = r^2 + v^2$.

$V = \frac{1}{3} \pi r^2 v$

$S = \pi r^2 + \pi r s$

Komolý kužel se dá představit jako kužel, kterému byla uříznuta špička rovinou rovnoběžnou s rovinou
podstavy, nebo také rotací pravoúhlého lichoběžníku. Komolý kužel má dvě podstavy s poloměry $r_1,
r_2$.

$V = \frac{1}{3} \pi v(r_1^2 + r_1 r_2 + r_2^2)$

$S = \pi r_1^2 + \pi r_2^2 + \pi s(r_1 + r_2)$
\subsubsection*{Koule a spol}
Koule vzniká rotací půlkruhu kolem jeho průměru. Kulová plocha neboli sféra je pouze povrch či slupka
koule, tedy nemá žádný objem ale pouze povrch, který je totožný s povrchem dané koule.

$V = \frac{4}{3} \pi r^3$

$S = 4 \pi r^2$

Polokoule je polovina koule.

$V = \frac{2}{3} \pi r^3$

$S = 2 \pi r^2 + \pi r^2 = 3 \pi r^2$

Kulový vrchlík je část kulové plochy, která je odseknuta rovinou. Dá se představit jako taková 
čepička. Stejně jako kulová plocha nemá žádný objem.

$S = 2 \pi r h$, $h$ je výška vrchlíku.

Kulová úseč je těleso, které vznikne odesknutím části koule jednou rovinou. Dá se říct, že vrchlík
je plášť kulové úseče. $q$ je poloměr podstavy kulové úseče a $h$ je výška úseče.

$V = \frac{1}{6} \pi h(3 q^2 + h^2) = \pi h^2(r - \frac{h}{3})$

$S = 2\pi r h + \pi q^2$

Pro kulovou úseč platí $r^2 = q^2 + (r - h)^2$

Představ si že vložíš do koule kužel s vrcholem ve středu koule, pokud zbytek koule vyřízneme,
získáme kulová výseč, která je podobná kornoutu se zmrzlinou.

$V = \frac{2}{3} \pi r^2 h$

$S = 2 \pi r h$

\subsection*{Cavalieriův princip}
Cavalieriův princip nám říká, že tělesa se stejně velkými podstavami a výškami mají stejný objem,
pokud mají řezy rovnoběžné s podstavani a vedené ve stejné vzdálenosti od podstav stejné obsahy.
Kvůli tomu mají kolmá a kosá tělesa stejný objem, pokud mají stejnou výšku a velkikost podstavy.
Stejně tak kužel a jehlan se stejnou výškou obsahem základny má také stejný objem.
\end{document}
